\chapter{Introduction}
\label{chap:introduction}

This blueprint describes the mathematics formalized in the Lean library
\href{https://github.com/glaroton/lc-manifolds}{lc-manifolds}: manifolds modeled on locally convex
spaces, their tangent bundles, vector fields and flows, differential forms, submanifolds and
immersions, vector bundles, connections and geodesics, metrics, manifolds of mappings, analytic,
complex and symplectic manifolds, and a collection of examples from functional analysis and from
groups of diffeomorphisms. Every definition and
theorem below is linked to its Lean counterpart, and the dependency graph shows how the results
build on each other. All the statements are proved in Lean, with no \texttt{sorry}, and the main
theorems depend only on the standard axioms of Lean.

\section*{The setting}

The differential calculus is the one of Michal and Bastiani, formalized in the companion library
lc-analysis \cite{LCA}: a map between open subsets of locally convex spaces is $\Ck{1}$ if its
directional derivatives exist and depend continuously on the point and on the direction, jointly.
This is the calculus of infinite-dimensional Lie theory, as presented by Neeb \cite{Nee06} and by
Glöckner and Neeb \cite{GN26}, which we follow for the foundations. For the geometry of Banach
manifolds (sprays, connections, Riemannian metrics) we follow Lang \cite{Lan99} and the book
\cite{Lar24}, and for metric geometry \cite{Lar19}.

A manifold modeled on $E$ is a charted space whose transition maps are smooth diffeomorphisms
between open subsets of $E$. Mathlib \cite{mC20} accepts an arbitrary topological model space in
its notion of charted space, so we do not need a separate theory of manifolds: we only add the
groupoid of Michal--Bastiani diffeomorphisms.

\section*{Conventions}

Three conventions are used throughout, and we state them once and for all.

\begin{itemize}
  \item \textbf{Model spaces need not be complete.} The model space $E$ is a real locally convex
  Hausdorff space, with no completeness assumption, following Glöckner and Neeb \cite{GN26}. We
  call a manifold \emph{Banach}, \emph{Hilbert} or \emph{Fréchet} only when its model is complete;
  a manifold modeled on a normed space that is not complete is not a Banach manifold. Whenever a
  result needs completeness, it is part of the statement.

  \item \textbf{No inverse function theorem beyond Banach spaces.} Outside Banach spaces there is
  no inverse function theorem, and there are no existence theorems for ordinary differential
  equations. On general models we avoid both: submanifolds are given by slice charts, examples are
  built with explicit charts, and geodesics are characterized, not constructed. Flows, geodesics,
  the exponential map and the Frobenius and Darboux theorems are proved for Banach models only.

  \item \textbf{Infinite-dimensional subtleties are flagged.} Notions that coincide in finite
  dimensions often split apart here: immersions and weak immersions, embedded and initial
  submanifolds, strong and weak metrics, continuity on the tangent bundle and continuity along
  curves. We keep a separate name for each notion, and each splitting is illustrated by a
  formalized counterexample whenever we have one.
\end{itemize}

\section*{Metrics: two questions, two words}

Chapter~\ref{chap:metrics} uses two pairs of words that should not be confused. For a field of
Finsler seminorms $\mu$, the question is how $\mu_p$ depends on $p$. A \emph{Finsler metric} is
continuous on the tangent bundle; this is what the term means without further qualification. A
\emph{Finsler length structure} is only continuous along $\Ck{1}$ curves: it is what one needs to
measure lengths, and quotient metrics on homogeneous spaces are of this kind. Every Finsler metric
is a Finsler length structure, and the converse fails already on $\R^2$. For a field of bilinear
forms $g$, the question is which topology $g_x$ induces on the tangent space: a \emph{strong}
metric induces the topology of the model, a \emph{weak} one may induce a coarser topology. Weak
Riemannian metrics are Finsler metrics; Lorentzian and pseudo-Riemannian metrics, of any index, are
not.

\section*{Organization}

The chapters follow the folders of the library.

\begin{itemize}
  \item \textbf{Foundations.} Manifolds and smooth maps (Chapter~\ref{chap:manifolds}), tangent
    bundles (Chapter~\ref{chap:tangent}), vector fields and the Frobenius theorem
    (Chapter~\ref{chap:vectorfields}), flows (Chapter~\ref{chap:flows}), manifolds with rough
    boundary (Chapter~\ref{chap:roughboundary}), and submanifolds and immersions
    (Chapter~\ref{chap:submanifolds}).
  \item \textbf{Forms and bundles.} Differential forms (Chapter~\ref{chap:forms}), vector bundles
    (Chapter~\ref{chap:bundles}), connections (Chapter~\ref{chap:connections}), and sprays and
    geodesics (Chapter~\ref{chap:geodesics}).
  \item \textbf{Metrics.} Metric geometry of length spaces (Chapter~\ref{chap:metric-spaces}), and
    Finsler, Riemannian, Lorentzian and pseudo-Riemannian metrics on manifolds
    (Chapter~\ref{chap:metrics}).
  \item \textbf{Further structures.} Local additions and manifolds of mappings
    (Chapter~\ref{chap:mapping}), analytic manifolds (Chapter~\ref{chap:analytic}), complex
    manifolds (Chapter~\ref{chap:complex}) and symplectic manifolds (Chapter~\ref{chap:symplectic}).
  \item \textbf{Examples} (Chapter~\ref{chap:examples}) and \textbf{counterexamples}
    (Chapter~\ref{chap:counterexamples}).
\end{itemize}

\section*{How the library was written}

@glaroton set the mathematical program, chose the framework and the references, fixed
the design constraints and reviewed the results. The Lean proofs were written by the AI assistant
Claude (Anthropic), working under his direction. Correctness does not rest on trusting either of
them: Lean's kernel \cite{dMU21} checks every proof, and the statements are kept close to the
literature so that they can be compared with it directly.
